\documentclass[11pt,a4paper]{article}

\usepackage[margin=2.5cm, headheight=15pt]{geometry}
\usepackage{amsmath,amssymb,amsthm}
\usepackage{mathtools}
\usepackage{microtype}
\usepackage{hyperref}
\usepackage{cleveref}
\usepackage{enumitem}
\usepackage{mdframed}
\usepackage{xcolor}
\usepackage{titlesec}
\usepackage{booktabs}
\usepackage{tabularx}
\usepackage{fancyhdr}

\definecolor{navyblue}{RGB}{26,58,92}
\definecolor{midblue}{RGB}{44,95,138}
\definecolor{lightblue}{RGB}{232,240,248}
\definecolor{verylightblue}{RGB}{248,251,254}
\definecolor{proofgray}{RGB}{242,242,242}

\hypersetup{colorlinks=true,linkcolor=midblue,citecolor=midblue,urlcolor=midblue,
  pdftitle={SSP Research Analysis},bookmarksnumbered=true}

\pagestyle{fancy}\fancyhf{}
\fancyhead[L]{\small\color{navyblue}\textit{SSP: Heuristic Bounds, Grouping Selection, Collapse Variants}}
\fancyhead[R]{\small\color{navyblue}\thepage}
\renewcommand{\headrulewidth}{0.4pt}
\renewcommand{\headrule}{\hbox to\headwidth{\color{navyblue}\leaders\hrule height\headrulewidth\hfill}}

\titleformat{\section}{\large\bfseries\color{navyblue}}{}{0em}{}[\color{navyblue}\titlerule]
\titleformat{\subsection}{\normalsize\bfseries\color{midblue}}{}{0em}{}
\titleformat{\subsubsection}{\normalsize\bfseries\itshape\color{midblue}}{}{0em}{}
\titlespacing{\section}{0pt}{18pt}{8pt}
\titlespacing{\subsection}{0pt}{12pt}{5pt}
\titlespacing{\subsubsection}{0pt}{8pt}{3pt}

\newmdenv[backgroundcolor=verylightblue,linecolor=navyblue,linewidth=1.2pt,
  leftline=true,rightline=false,topline=false,bottomline=false,
  innerleftmargin=12pt,innerrightmargin=8pt,innertopmargin=8pt,innerbottommargin=8pt,
  skipabove=10pt,skipbelow=6pt]{thmbox}

\newmdenv[backgroundcolor=proofgray,linecolor=midblue,linewidth=0.8pt,
  leftline=true,rightline=false,topline=false,bottomline=false,
  innerleftmargin=10pt,innerrightmargin=8pt,innertopmargin=6pt,innerbottommargin=6pt,
  skipabove=6pt,skipbelow=6pt]{proofbox}

\newmdenv[backgroundcolor=lightblue,linecolor=navyblue,linewidth=1.5pt,roundcorner=3pt,
  innerleftmargin=14pt,innerrightmargin=14pt,innertopmargin=10pt,innerbottommargin=10pt,
  skipabove=12pt,skipbelow=8pt]{constructbox}

\theoremstyle{plain}
\newtheorem{theorem}{Theorem}[section]
\newtheorem{lemma}[theorem]{Lemma}
\newtheorem{corollary}[theorem]{Corollary}
\newtheorem{proposition}[theorem]{Proposition}
\theoremstyle{definition}
\newtheorem{definition}[theorem]{Definition}
\newtheorem{construction}[theorem]{Construction}
\newtheorem{remark}[theorem]{Remark}

\DeclareMathOperator{\OPT}{OPT}
\DeclareMathOperator*{\argmin}{argmin}
\DeclareMathOperator*{\argmax}{argmax}
\newcommand{\Cconf}{\mathfrak{C}}
\newcommand{\calG}{\mathcal{G}}
\newcommand{\Kstar}{K^{*}}
\newcommand{\sw}[2]{w(#1,#2)}

\setlist[itemize]{noitemsep,topsep=4pt}
\setlist[enumerate]{noitemsep,topsep=4pt}

% ─────────────────────────────────────────────────────────────────────────────
\begin{document}
% ─────────────────────────────────────────────────────────────────────────────

\begin{titlepage}
\centering\vspace*{2cm}
{\color{navyblue}\rule{\textwidth}{2pt}}\\[0.5cm]
{\Huge\bfseries\color{navyblue}Research Analysis}\\[0.4cm]
{\LARGE\bfseries\color{navyblue}Job Sequencing and Tool Switching Problem}\\[0.3cm]
{\Large\itshape\color{midblue}Heuristic Bounds, Grouping Selection, and Collapse Variants}\\[0.3cm]
{\color{navyblue}\rule{\textwidth}{2pt}}\\[1.5cm]
{\normalsize\begin{tabular}{c}
Part~I\quad Worst-Case Gap: Configuration-Based TSP Heuristic vs.\ OPT-SSP\\[2pt]
Part~II\quad Grouping Selection: Mathematical and ML Directions\\[2pt]
Part~III\quad SSP Variants that Collapse to JGP
\end{tabular}}\\[2cm]
\begin{mdframed}[backgroundcolor=lightblue,linecolor=navyblue,linewidth=1pt,
  innerleftmargin=18pt,innerrightmargin=18pt,innertopmargin=12pt,innerbottommargin=12pt]
\small
\textbf{Preface on novelty.}\ Several foundational results used here are from the
literature and are clearly attributed: the KTNS heuristic and the SSP formulation
\cite{TangDenardo1988}; NP-hardness of SSP \cite{Crama1994}; the LP lower bound
$\OPT_{\mathrm{SSP}} \ge m-C$ and column generation for SSP \cite{daSilva2021};
the Hamiltonian-cycle reformulation of SSP \cite{Ghiani2007}.
The \emph{original contributions} of this document are:
(i)~a tight worst-case bound on the padding-choice gap of the canonical-configuration
TSP heuristic with a complete instance construction and proof
(\cref{sec:partI});
(ii)~the equivalence between optimal group ordering and a Max-Weight Hamiltonian Path
on the tool-overlap graph (\cref{thm:overlap}); and
(iii)~the general characterisation of cost functions under which SSP collapses to JGP,
together with novel structured variants (\cref{sec:partIII}).
\end{mdframed}
\vfill
\end{titlepage}

\tableofcontents\newpage

% ─────────────────────────────────────────────────────────────────────────────
\section{Notation and Background}
\label{sec:notation}
% ─────────────────────────────────────────────────────────────────────────────

Throughout: $J=\{1,\ldots,n\}$ jobs; $T=\{1,\ldots,m\}$ tools; $C$ magazine capacity;
$T_j\subseteq T$, $|T_j|\le C$: tool requirement of job $j$.

A \textbf{configuration} for group $G_k$ is any $\mathcal{C}_k\subseteq T$ with
$|\mathcal{C}_k|=C$ and $\bigcup_{j\in G_k}T_j\subseteq\mathcal{C}_k$.
$\Cconf(G_k)$: all valid configurations for $G_k$.
Switching cost $\sw{\mathcal{C}}{\mathcal{C}'}=|\mathcal{C}\setminus\mathcal{C}'|$
(tools removed $=$ tools inserted under the uniform model \cite{TangDenardo1988}).

A \textbf{grouping} $\calG=\{G_1,\ldots,G_K\}$ is a partition of $J$ into $K$ groups.
$\Kstar$: JGP-optimal (minimum) number of groups.
$s_k = C-|T_{G_k}|$: padding slack of group $G_k$.

$\OPT_{\mathrm{SSP}}$: optimal SSP cost, minimised over \emph{all} groupings $\calG$,
all valid configurations $\mathcal{C}_k\in\Cconf(G_k)$, and all orderings of groups
and jobs within groups.

\medskip
\noindent\textbf{The heuristic under study.}
(i)~Solve JGP to obtain $\Kstar$ groups.
(ii)~Assign each group $G_k$ a \emph{canonical} configuration
$\hat{\mathcal{C}}_k\in\Cconf(G_k)$: padding slots are filled without reference
to adjacent groups' mandatory tools (e.g.\ all padding filled by a fixed shared
filler set, or independently at random from remaining tools).
(iii)~Solve TSP on the $\Kstar$ configurations with arc weight $\sw{\cdot}{\cdot}$
to find the best ordering.
(iv)~Read off the SSP path. Cost: $H$.

\medskip
\noindent\textbf{Where the gap comes from.}
The TSP in step~(iii) finds the optimal \emph{ordering} of the given configurations;
it does not alter the configurations themselves.
The gap $H-\OPT_{\mathrm{SSP}}\ge 0$ therefore has two sources:
\begin{enumerate}[label=(\arabic*)]
  \item \emph{Padding-choice gap}: the canonical configuration for $G_k$ may waste
        padding slots on tools that are irrelevant to adjacent groups, forcing more
        switches than necessary. Choosing padding from adjacent groups' mandatory
        tools (``lookahead'') reduces switches.
  \item \emph{Grouping-choice gap}: OPT-SSP may use a grouping with $K>\Kstar$ groups,
        where the extra configuration boundaries allow tool layouts that reduce total
        switches below anything achievable with $\Kstar$ groups.
\end{enumerate}
Part~I isolates and bounds Source~(1) for a fixed JGP-optimal grouping, since
that is the gap introduced specifically by the heuristic's padding choice.
Source~(2) is a property of the JGP/SSP relationship itself and is not attributable
to the heuristic; it is discussed separately in \cref{rem:source2}.

% ─────────────────────────────────────────────────────────────────────────────
\section{Part~I: Tight Bound on the Padding-Choice Gap}
\label{sec:partI}
% ─────────────────────────────────────────────────────────────────────────────

\subsection{Setup and Key Lemma}

For a fixed JGP-optimal grouping $\calG^*=\{G_1,\ldots,G_{\Kstar}\}$ and a fixed
ordering $\pi$ of the groups, the switching cost at each boundary depends on how
much of $T_{G_{\pi(k)}}$ survives into $\mathcal{C}_{\pi(k+1)}$.
The mandatory tools of $G_{\pi(k)}$ that are \emph{not} required by $G_{\pi(k+1)}$
will survive only if they are placed in the $s_{\pi(k+1)}=C-|T_{G_{\pi(k+1)}}|$
padding slots of $\mathcal{C}_{\pi(k+1)}$.

\begin{thmbox}
\begin{lemma}[Achievable Per-Transition Cost {\normalfont (new)}]
\label{lem:transition}
For consecutive groups $G_k$, $G_{k+1}$ in a fixed ordering, the minimum switching
cost achievable by any choice of $\mathcal{C}_k\in\Cconf(G_k)$ and
$\mathcal{C}_{k+1}\in\Cconf(G_{k+1})$ is
\[
  w^*_k
  \;=\;
  \max\!\bigl(0,\;|T_{G_k}\setminus T_{G_{k+1}}|-s_{k+1}\bigr)
  \;=\;
  \max\!\bigl(0,\;|T_{G_k}\cup T_{G_{k+1}}|-C\bigr),
\]
and it is achieved by the \emph{lookahead} configuration
$\mathcal{C}_k^*=T_{G_k}\cup X_k$ where $X_k\subseteq T_{G_{k+1}}\setminus T_{G_k}$,
$|X_k|=\min(s_k,\,|T_{G_{k+1}}\setminus T_{G_k}|)$.
\end{lemma}
\end{thmbox}

\begin{proofbox}
\begin{proof}
$\mathcal{C}_{k+1}$ must allocate $|T_{G_{k+1}}|$ slots to mandatory tools,
leaving $s_{k+1}$ padding slots.
Of the $|T_{G_k}\setminus T_{G_{k+1}}|$ tools in $\mathcal{C}_k$ that $G_{k+1}$
does not need, at most $s_{k+1}$ can be kept in $\mathcal{C}_{k+1}$ as padding.
The remaining $\max(0,|T_{G_k}\setminus T_{G_{k+1}}|-s_{k+1})$ must leave,
giving the lower bound. Equality: set $X_k$ to be the first $\min(s_k,r-ov)$ tools
of $T_{G_{k+1}}\setminus T_{G_k}$ where $r=|T_{G_k}|$; these survive into
$\mathcal{C}_{k+1}^*$ since $X_k\subseteq T_{G_{k+1}}$.
Direct computation gives $|\mathcal{C}_k^*\setminus\mathcal{C}_{k+1}^*|=w^*_k$. \qed
\end{proof}
\end{proofbox}

\begin{remark}
\label{rem:literature-lb}
The global lower bound $\OPT_{\mathrm{SSP}}\ge m-C$ due to~\cite{daSilva2021}
(every tool beyond the initial magazine load must be inserted at least once)
is a bound on the SSP over all groupings.
\cref{lem:transition} is a \emph{per-transition} bound for a fixed grouping and
ordering; it is a different and complementary tool.
\end{remark}

\subsection{The Tight Construction}
\label{sec:construction}

\begin{constructbox}
\begin{construction}[Instance $\mathcal{I}(K,r,\mathit{ov},C)$]
\label{const:main}
\textbf{Parameters.} Positive integers satisfying
\begin{equation}
  \label{eq:params}
  0\le \mathit{ov}\le\tfrac{r}{2},
  \quad
  r < C,
  \quad
  2r-\mathit{ov} > C.
\end{equation}
Consistency: $\mathit{ov}\le r/2$ and $2r-\mathit{ov}>C>r$ require $r>C/2$;
e.g.\ $r=6,\,\mathit{ov}=2,\,C=9$ satisfies $2(6)-2=10>9>6$ and $2\le 3$. \medskip

\noindent\textbf{Tool universe.}
$T=\{t_1,\ldots,t_m\}$, $m=(K-1)(r-\mathit{ov})+r$.\medskip

\noindent\textbf{Mandatory tool sets.} For $k=1,\ldots,K$:
\[
  T_{G_k}=\{t_{(k-1)(r-\mathit{ov})+1},\ldots,t_{(k-1)(r-\mathit{ov})+r}\},
\]
consecutive windows of width $r$ shifted by $r-\mathit{ov}$.
Thus $|T_{G_k}|=r$,\ $|T_{G_k}\cap T_{G_{k+1}}|=\mathit{ov}$,\
$|T_{G_k}\cup T_{G_{k+1}}|=2r-\mathit{ov}>C$.\medskip

\noindent\textbf{Jobs.} $G_k$ contains one job per tool subset of $T_{G_k}$
(or simply one job requiring all of $T_{G_k}$); the mandatory tool set of the group
is $T_{G_k}$. Padding slack $s=C-r>0$ for all groups.\medskip

\noindent\textbf{JGP-optimality ($\Kstar=K$).}
Since $|T_{G_k}\cup T_{G_{k+1}}|>C$ for all $k$, no configuration covers two adjacent
groups, so $\Kstar=K$.
\end{construction}
\end{constructbox}

\subsubsection{Heuristic Cost}

Use canonical configuration $\hat{\mathcal{C}}_k=T_{G_k}\cup F$ where
$F\subseteq T\setminus\bigcup_k T_{G_k}$ is a fixed filler set of $s$ tools
disjoint from all mandatory sets.
(If $m=K r-(K-1)\mathit{ov}$ leaves no room for such $F$, use
$F_k\subseteq T\setminus(T_{G_{k-1}}\cup T_{G_k}\cup T_{G_{k+1}})$;
for large enough $K$ this is non-empty.)
Then for consecutive groups in the ordering:
\begin{align*}
  \hat{\mathcal{C}}_k\setminus\hat{\mathcal{C}}_{k+1}
  &=\underbrace{(T_{G_k}\setminus T_{G_{k+1}})}_{r-\mathit{ov}\text{ tools}}
   \;\cup\;
   \underbrace{(F\setminus\hat{\mathcal{C}}_{k+1})}_{F\cap T_{G_{k+1}}=\emptyset,\text{ so all }s\text{ tools leave}},
\end{align*}
giving $\sw{\hat{\mathcal{C}}_k}{\hat{\mathcal{C}}_{k+1}}=
(r-\mathit{ov})+s = C-\mathit{ov}$ per transition. The TSP orders the $K$
configurations optimally; since all pairwise costs are symmetric and equal by the
window structure, any Hamiltonian path has the same cost:
\begin{equation}
  \label{eq:H}
  H = (K-1)(C-\mathit{ov}).
\end{equation}

\subsubsection{Optimal Cost for This Grouping Under Lookahead Configurations}
\label{sec:lookahead}

For the \emph{same} JGP-optimal $K$-group partition, use lookahead configuration
$\mathcal{C}_k^*=T_{G_k}\cup X_k$ with
$X_k\subseteq T_{G_{k+1}}\setminus T_{G_k}$, $|X_k|=s=C-r$.
This is valid since $|T_{G_{k+1}}\setminus T_{G_k}|=r-\mathit{ov}>s$
(by $2r-\mathit{ov}>C$, i.e.\ $r-\mathit{ov}>C-r=s$).

\medskip\noindent Transition cost in the natural ordering
$G_1\to G_2\to\cdots\to G_K$:
\begin{align*}
  \mathcal{C}_k^*\setminus\mathcal{C}_{k+1}^*
  &= (T_{G_k}\setminus\mathcal{C}_{k+1}^*)\cup(\underbrace{X_k\setminus\mathcal{C}_{k+1}^*}_{=\,\emptyset,\;X_k\subseteq T_{G_{k+1}}}).
\end{align*}
Now $\mathcal{C}_{k+1}^*=T_{G_{k+1}}\cup X_{k+1}$ with
$X_{k+1}\subseteq T_{G_{k+2}}\setminus T_{G_{k+1}}$.
Since $\mathit{ov}\le r/2$, windows at distance~2 are disjoint:
$T_{G_k}\cap T_{G_{k+2}}=\emptyset$.
Hence $T_{G_k}\cap\mathcal{C}_{k+1}^*=T_{G_k}\cap T_{G_{k+1}}=\mathit{ov}$ tools, so
\[
  |\mathcal{C}_k^*\setminus\mathcal{C}_{k+1}^*|=r-\mathit{ov}.
\]
Total cost of the lookahead solution on this grouping:
\begin{equation}
  \label{eq:lookahead}
  \mathrm{Cost}_{\text{lookahead}}=(K-1)(r-\mathit{ov}).
\end{equation}
This is an \emph{upper bound} on $\OPT_{\mathrm{SSP}}$:
\begin{equation}
  \label{eq:UB}
  \OPT_{\mathrm{SSP}}\le (K-1)(r-\mathit{ov}).
\end{equation}

\subsection{Lower Bound on \texorpdfstring{$\OPT_{\mathrm{SSP}}$}{OPT-SSP} for This Instance}
\label{sec:LB}

We now prove the matching lower bound.

\begin{thmbox}
\begin{lemma}[Lower Bound via Mandatory Tool Exits]
\label{lem:LB}
For instance $\mathcal{I}(K,r,\mathit{ov},C)$, any feasible SSP solution---over
all possible groupings, all valid configurations, and all job orderings---incurs
at least $(K-1)(r-\mathit{ov})$ tool switches.
\end{lemma}
\end{thmbox}

\begin{proofbox}
\begin{proof}
\textbf{Define the mandatory exit sets.}
For $k=1,\ldots,K-1$ let
\[
  D_k = T_{G_k}\setminus T_{G_{k+1}}
      = \{t_{(k-1)(r-\mathit{ov})+1},\ldots,t_{(k-1)(r-\mathit{ov})+(r-\mathit{ov})}\},
\]
the $r-\mathit{ov}$ tools required by jobs in $G_k$ but by no job in
$G_{k+1},G_{k+2},\ldots,G_K$.
This last claim follows from the window structure: tools in $D_k$ lie at positions
$(k-1)(r-\mathit{ov})+1$ through $k(r-\mathit{ov})$ in the global tool list,
strictly to the left of all windows $G_{k+1},\ldots,G_K$.
The sets $D_1,\ldots,D_{K-1}$ are pairwise disjoint by construction.

\medskip
\textbf{Each tool in $D_k$ causes at least one switch in any feasible solution.}
Fix any tool $t\in D_k$ and any feasible SSP solution (any grouping, any job sequence).
\begin{itemize}
  \item $t$ is required by at least one job $j\in G_k$. Hence $t$ must be present
        in the magazine when $j$ is processed. Since the magazine starts with some
        initial load of $C$ tools, $t$ is either in the initial load (zero insertions
        so far) or must be inserted at some point before $j$ is processed (one insertion).
  \item After the last job requiring $t$ is processed (and $t\in D_k$ means no job
        outside $G_k$ requires $t$), tool $t$ is never needed again. Because the
        magazine has finite capacity $C$ and the SSP continues processing further jobs
        that may require tools not currently loaded, $t$ will be displaced from the
        magazine at some point. That displacement is one \emph{removal}.
  \item In the uniform SSP model~\cite{TangDenardo1988}, the magazine stays at
        capacity $C$ throughout: every removal is paired with an insertion and vice
        versa. So either $t$ was in the initial load and gets removed once (one switch),
        or $t$ was not in the initial load, gets inserted once and removed once (two
        switches). In both cases, $t$ contributes at least \textbf{one switch}
        (the removal) to the total switch count.
\end{itemize}

\medskip
\textbf{This argument is independent of the grouping used.}
No matter how the solution groups and sequences jobs, the requirement that $t\in D_k$
must be present when some job in $G_k$ runs, and absent (eventually) afterwards,
is a property of the \emph{instance} (which jobs require which tools), not of the
solution structure. A finer grouping (more configurations) does not change which tools
are required by which jobs.

\medskip
\textbf{Count.}
Since $D_1,\ldots,D_{K-1}$ are pairwise disjoint, the $\sum_{k=1}^{K-1}|D_k|
=(K-1)(r-\mathit{ov})$ removal events for tools in $\bigcup_k D_k$ are all distinct
switch events. Hence total switches $\ge (K-1)(r-\mathit{ov})$. \qed
\end{proof}
\end{proofbox}

\begin{remark}[On Source~(2) of the gap]
\label{rem:source2}
\cref{lem:LB} shows that $\OPT_{\mathrm{SSP}}=(K-1)(r-\mathit{ov})$ for this
instance: the lower bound matches the lookahead upper bound~\eqref{eq:UB}.
Importantly, the lower bound holds for any grouping, including those with $K'>K$
groups: a finer grouping cannot reduce the mandatory removal count of tools in
$\bigcup_k D_k$, because tool requirements are fixed by the instance.
Hence for $\mathcal{I}(K,r,\mathit{ov},C)$, Source~(2) of the gap
(grouping choice) contributes \emph{zero}: the optimal SSP uses exactly $K=\Kstar$
groups, and the entire gap is Source~(1) (padding choice).
This makes the construction a clean isolator of Source~(1).
\end{remark}

\subsection{Main Result of Part~I}

\begin{thmbox}
\begin{theorem}[Tight Padding-Choice Gap]
\label{thm:gap}
For instance $\mathcal{I}(K,r,\mathit{ov},C)$ with parameters satisfying~\eqref{eq:params}:
\begin{align*}
  H &= (K-1)(C-\mathit{ov}), \\
  \OPT_{\mathrm{SSP}} &= (K-1)(r-\mathit{ov}), \\
  H - \OPT_{\mathrm{SSP}} &= (K-1)(C-r) = (K-1)\,s,
\end{align*}
where $s=C-r$ is the per-group padding slack.
The bound is \textbf{tight}: it is achieved by the construction above.
No canonical-padding TSP heuristic can guarantee a smaller gap on this family.
\end{theorem}
\end{thmbox}

\begin{proof}
$H$ follows from~\eqref{eq:H}. $\OPT_{\mathrm{SSP}}=(K-1)(r-\mathit{ov})$
follows from the upper bound~\eqref{eq:UB} and \cref{lem:LB}.
The gap $=(K-1)(C-\mathit{ov})-(K-1)(r-\mathit{ov})=(K-1)(C-r)$.
Tightness: the heuristic cost~\eqref{eq:H} is a worst-case cost for canonical padding
(any filler $F$ disjoint from adjacent mandatory sets gives the same cost);
no canonical strategy does better since it cannot know adjacency. \qed
\end{proof}

\begin{corollary}[Unbounded Approximation Ratio]
\label{cor:ratio}
\[
  \frac{H}{\OPT_{\mathrm{SSP}}}=\frac{C-\mathit{ov}}{r-\mathit{ov}}=1+\frac{s}{r-\mathit{ov}}.
\]
As $s\to r-\mathit{ov}-1$ (slack approaches maximum under the constraints), the ratio
approaches $C-\mathit{ov}$, which grows with $C$.
As $r-\mathit{ov}\to 1$, the ratio approaches $C-\mathit{ov}+1\approx C$.
The ratio is unbounded as $C\to\infty$: \textbf{no constant approximation factor}
exists for the canonical-padding TSP heuristic.
\end{corollary}

\begin{remark}[Gap is zero when $s=0$]
When $r=C$ (jobs fill the magazine completely), no padding freedom exists and every
configuration is forced. SSP then reduces to a pure TSP on mandatory tool sets,
consistent with the known case where $|T_j|=C$ for all $j$~\cite{Ghiani2007}.
\end{remark}

% ─────────────────────────────────────────────────────────────────────────────
\section{Part~II: Grouping Selection}
\label{sec:partII}
% ─────────────────────────────────────────────────────────────────────────────

\noindent\textbf{Context.}
You have: all JGP-optimal groupings ($K=\Kstar$), some sub-optimal ones ($K>\Kstar$),
and the trivial all-configurations option.
The goal: select groupings to (a) solve SSP exactly or (b) obtain strong heuristic
solutions with quality guarantees.

\subsection{Mathematical Directions}

\subsubsection{GTSP Formulation (Reference and Extension)}
\label{sec:GTSP}

\noindent\textbf{Known.}
Ghiani et al.~\cite{Ghiani2007} showed that SSP is equivalent to a minimum-cost
Hamiltonian cycle on a graph of all valid configurations.
Column generation for SSP with set-covering master and bin-packing pricing was
developed in~\cite{daSilva2021}.

\medskip
\noindent\textbf{Extension to fixed grouping (new framing).}
For a \emph{fixed} grouping $\calG=\{G_1,\ldots,G_K\}$, the subproblem of jointly
selecting configurations and ordering groups is a Generalised TSP (GTSP):
clusters are $V_k=\Cconf(G_k)$, arc costs are $\sw{\mathcal{C}_i}{\mathcal{C}_j}$.
This is the natural decomposition when iterating over candidate groupings.

\begin{thmbox}
\begin{theorem}[GTSP Equivalence for Fixed Grouping]
\label{thm:GTSP}
For a fixed grouping $\calG$, $\OPT(\calG)=$ minimum-cost Hamiltonian path
in the GTSP on clusters $\{V_k\}_{k=1}^K$.
\end{theorem}
\end{thmbox}

\begin{proof}
Any feasible SSP path on $\calG$ selects one configuration per group and an ordering,
corresponding to exactly a Hamiltonian path visiting one node per cluster.
Minimising over both choices is the GTSP. \qed
\end{proof}

\noindent\textbf{Solving GTSP.}
The Noon--Bean transformation~\cite{NoonBean1993} reduces GTSP to standard TSP
on $O(N)$ nodes ($N=\sum_k|V_k|$); Concorde or LKH3 apply directly.
Held--Karp DP for GTSP runs in $O(N^2\cdot 2^K)$, tractable for small $\Kstar$.

\noindent\textbf{Cluster size.}
$|V_k|=\binom{m-|T_{G_k}|}{s_k}$ is exponential in $s_k$.
Restrict to \emph{structured configurations}:
$\mathcal{C}_k=T_{G_k}\cup X_k$ with $X_k\subseteq T_{G_{k+1}}\setminus T_{G_k}$
(lookahead, \cref{sec:construction}); this limits each cluster to $O(r-\mathit{ov})$
nodes and recovers the tight construction of Part~I as the GTSP optimum.

\subsubsection{Optimal Ordering via Tool-Overlap Graph (New Result)}
\label{sec:overlap}

\begin{definition}[Tool-Overlap Graph]
For grouping $\calG=\{G_1,\ldots,G_K\}$, define the complete weighted graph
$H_{\mathrm{ov}}$ on $K$ nodes with edge weight
$w(G_k,G_{k'})=|T_{G_k}\cap T_{G_{k'}|}$.
\end{definition}

\begin{thmbox}
\begin{theorem}[Optimal Ordering = Max-Weight Hamiltonian Path (New)]
\label{thm:overlap}
For a fixed grouping with $|T_{G_k}|=r$ for all $k$ and padding slack
$s\ge r-\min_{k}|T_{G_k}\setminus T_{G_{k+1}}|$ (sufficient to enable full
lookahead), the minimum SSP cost over all orderings is:
\[
  \OPT(\calG,\pi^*)
  =(K-1)r - \max_\pi\sum_{k=1}^{K-1}|T_{G_{\pi(k)}}\cap T_{G_{\pi(k+1)}}|.
\]
Hence the optimal ordering $\pi^*$ solves a Max-Weight Hamiltonian Path (MWHP)
on $H_{\mathrm{ov}}$, solvable in $O(K^2\cdot 2^K)$ by Held--Karp DP.
\end{theorem}
\end{thmbox}

\begin{proof}
Under lookahead configuration $\mathcal{C}_{\pi(k)}^*=T_{G_{\pi(k)}}\cup X_{\pi(k)}$
with $X_{\pi(k)}\subseteq T_{G_{\pi(k+1)}}\setminus T_{G_{\pi(k)}}$, $|X_{\pi(k)}|=s$,
all padding tools survive the transition (they belong to $T_{G_{\pi(k+1)}}$),
and the mandatory overlap $T_{G_{\pi(k)}}\cap T_{G_{\pi(k+1)}}$ also survives.
Thus
$|\mathcal{C}_{\pi(k)}^*\setminus\mathcal{C}_{\pi(k+1)}^*|
=|T_{G_{\pi(k)}}|-|T_{G_{\pi(k)}}\cap T_{G_{\pi(k+1)}}|
=r-w(G_{\pi(k)},G_{\pi(k+1)})$.
Summing over $k=1,\ldots,K-1$ and using $\sum_k r=(K-1)r=\text{const}$,
minimising total cost is equivalent to maximising $\sum_k w(G_{\pi(k)},G_{\pi(k+1)})$. \qed
\end{proof}

\noindent\textbf{Practical use.}
Enumerate candidate groupings; for each, build $H_{\mathrm{ov}}$ and solve MWHP
(tractable for small $\Kstar$); score groupings by MWHP value; solve GTSP exactly
for the top-$p$ groupings; return the best.
The MWHP score is a principled, polynomial-time proxy for $\OPT(\calG)$
under the lookahead assumption.

\subsubsection{Lagrangian Relaxation for the JGP/SSP Tradeoff}
\label{sec:lagrangian}

The JGP/SSP tension is captured by
\[
  L(\lambda)=\min_{\calG,\pi,\text{configs}}
  \bigl[\mathrm{cost}(\calG,\pi)+\lambda(|\calG|-\Kstar)\bigr],\quad\lambda\ge 0.
\]
By weak Lagrangian duality~\cite[Ch.~10]{Wolsey1998}:
$L(\lambda)\le\OPT_{\mathrm{SSP}}$ for all $\lambda\ge 0$.
The subgradient algorithm (standard; see~\cite{Wolsey1998}) traces the Pareto frontier
of (number of groups, switch cost), providing sub-optimal JGP groupings ($K>\Kstar$)
that may achieve lower switching cost than any $\Kstar$-group solution.
The tightest bound is $L(\lambda^*)=\max_{\lambda\ge 0}L(\lambda)$; the gap
$\OPT_{\mathrm{SSP}}-L(\lambda^*)$ equals the LP integrality gap for the joint problem.

\subsubsection{Column Generation (Reference + Grouping-Selection Use)}
\label{sec:CG}

Column generation for SSP was developed in~\cite{daSilva2021} with a set-covering
master over feasible groups and a bin-packing pricing subproblem.
We use it here in a complementary way: rather than solving SSP from scratch, we use
the LP dual variables $u_j$ from an intermediate CG iteration as a \emph{scoring
function} for candidate groupings.
A grouping $\calG$ with high $\sum_{G_k\in\calG}\sum_{j\in G_k}u_j$
is ``dual-cheap'' and should be prioritised for exact GTSP evaluation.
This provides a data-driven ranking of candidate groupings compatible with the
CG framework, without requiring a full branch-and-price run.

\subsubsection{Tool-Conflict Graph and Chromatic Structure}
\label{sec:conflict}

\begin{definition}[Tool Conflict Graph {\normalfont(standard; see~\cite{Crama1994})}]
$G_c=(J,E_c)$ where $(j,j')\in E_c\iff|T_j\cup T_{j'}|>C$.
\end{definition}

Valid groupings correspond to proper $\Kstar$-colourings of $G_c$ (each colour class
is a feasible group). Among all $\Kstar$-colourings, the MWHP criterion of
\cref{thm:overlap} selects the one whose colour classes have maximum pairwise mandatory
tool overlap in the optimal ordering.
This connects the chromatic structure of $G_c$ to the SSP objective
in a way not, to our knowledge, made explicit in prior work.

\subsection{Machine Learning Directions}

\subsubsection{GNN Grouping Scorer}
\label{sec:GNN}

\paragraph{Goal.} Score a large candidate pool without solving GTSP for each grouping.

\paragraph{Architecture.}
Encode the instance as a bipartite job--tool graph $B=(J\cup T,E)$
with $(j,t)\in E\iff t\in T_j$.
A two-layer Graph Attention Network (GAT) or GraphSAGE produces job embeddings
$h_j\in\mathbb{R}^d$.
For grouping $\calG$, represent group $G_k$ as
$\mathbf{g}_k=\frac{1}{|G_k|}\sum_{j\in G_k}h_j$.
A second GNN on the $K$-node grouping graph (edge features $|T_{G_k}\cap T_{G_{k'}}|$)
predicts $\OPT(\calG)$.

\paragraph{Justification.}
$\OPT(\calG)$ depends on the tool-overlap structure of the grouping, captured by
the $K$-node graph. GNNs are permutation-equivariant (the SSP cost is
permutation-invariant in groups) and are universal approximators on graphs~\cite{Xu2019};
hence the architecture is principled.

\paragraph{Training.} Solve instances up to moderate scale exactly; include both
$K=\Kstar$ and $K>\Kstar$ groupings so the model learns across the Pareto frontier.
Loss: $\mathbb{E}[|\hat{c}(\calG)-\OPT(\calG)|^2]$.

\paragraph{Inference.} Score all candidates in $O(K^2)$ per grouping; solve GTSP
for top-$p$; return best.

\paragraph{Guarantee.} If the GNN approximates $\OPT(\calG)$ within error $\varepsilon$
w.p.\ $\ge 1-\delta$, then the returned solution has cost
$\le\OPT_{\mathrm{SSP}}+\varepsilon+\eta$, where $\eta\to 0$ as $p\to\infty$.

\subsubsection{Reinforcement Learning for Joint Grouping and Sequencing}
\label{sec:RL}

RL naturally explores $K>\Kstar$, learning when more groups reduce total switches
without being constrained to any fixed grouping size.

\paragraph{MDP.} State: partial job-to-group assignment. Action: assign next job to
an existing feasible group or open a new one (masked to maintain $|T_{G_k}|\le C$).
Terminal reward $R_T=-\mathrm{cost}(\text{final grouping, GTSP-optimal configs})$.

\paragraph{Training.} REINFORCE with greedy rollout baseline~\cite{Kool2019}:
\[
  \nabla_\theta\mathcal{L}=-\mathbb{E}[(R_T-b)\nabla_\theta\log\pi_\theta].
\]
Architecture: Transformer over job feature vectors $[|T_j|,\phi_j]$
($\phi_j\in\{0,1\}^m$ the tool-requirement vector) following~\cite{Kool2019}.

\paragraph{Distinction from JGP solvers.} A pure JGP solver minimises $K$;
the RL policy minimises switching cost, discovering that $K>\Kstar$ can be beneficial.
The reward directly measures switching cost, not group count.

\paragraph{Guarantee.} Distributional:
$\mathbb{E}_{I\sim\mathcal{D}}[\mathrm{cost}(\pi_\theta(I))]\to
\mathbb{E}_{I\sim\mathcal{D}}[\OPT_{\mathrm{SSP}}(I)]$ under standard RL convergence
conditions. No worst-case bound.

\subsubsection{Learning-Augmented Column Generation}
\label{sec:LACG}

\paragraph{Setup.} CG for SSP~\cite{daSilva2021} is exact but the pricing problem
is costly. Train a predictor $P_\phi:(u,B)\to G'$ (dual variables $u$ and instance
graph $B$ to a candidate feasible group $G'$) on $(u^*,G^*)$ pairs from solved
instances.

\paragraph{Protocol.} At each CG iteration: query $P_\phi$; verify reduced cost
$c_{G'}-\sum_{j\in G'}u_j$ (cheap); if negative, add to restricted master;
else fall back to exact pricing.

\paragraph{Correctness.} Every ML-suggested column is verified before use;
exact LP optimality is certified by the fallback solver. Exact optimality is preserved;
ML only reduces the number of exact pricing calls. Analogous
approaches reduce pricing iterations by 50--80\% on related problems~\cite{Gasse2019}.

\subsubsection{Contrastive Siamese GNN for Grouping Ranking}
\label{sec:contrastive}

Train a Siamese GNN on pairs $(\calG_A,\calG_B)$ labelled by which has lower
$\OPT$. Loss: binary cross-entropy. At inference: tournament over $N$ candidates
in $O(N\log N)$ comparisons. Only the top-ranked grouping is passed to GTSP.
Requires only \emph{relative} quality labels, cheaper to obtain than exact costs.

% ─────────────────────────────────────────────────────────────────────────────
\section{Part~III: SSP Variants That Collapse to JGP}
\label{sec:partIII}
% ─────────────────────────────────────────────────────────────────────────────

\begin{definition}[Collapse to JGP]
\label{def:collapse}
SSP-$f$ (SSP with cost function $f:2^T\times 2^T\to\mathbb{R}_+$)
\emph{collapses to JGP} if for every instance $I$,
$\OPT_{\mathrm{SSP}\text{-}f}$ is achieved by some grouping with exactly $\Kstar(I)$ groups.
\end{definition}

\subsection{Baseline: Switching Instants (Known)}

$f_{01}(\mathcal{C},\mathcal{C}')=\mathbf{1}[\mathcal{C}\ne\mathcal{C}']$: cost equals
number of configuration changes $=K-1$ for $K$ groups. Minimising $\Rightarrow$ JGP.
This is the Tool Switching Instants Problem (TSIP), known to be equivalent to
JGP~\cite{TangDenardo1988}. It serves as the baseline; we seek metrically richer variants.

\subsection{General Collapse Characterisation (New)}
\label{sec:general}

Rather than analysing specific variants individually, we first characterise all
functions $f$ that cause collapse. This is new and subsumes the threshold and
weighted variants as corollaries.

\begin{thmbox}
\begin{theorem}[Necessary and Sufficient Conditions for Collapse]
\label{thm:collapse}
$f$ causes SSP-$f$ to collapse to JGP on all instances if and only if:
\begin{description}
  \item[(C1) Self-cost zero:] $f(\mathcal{C},\mathcal{C})=0$ for all $\mathcal{C}$.
  \item[(C2) Strict inter-group positivity:] For any JGP-optimal grouping, any two
        groups $G_k\ne G_{k'}$ in it, and any valid configurations
        $\mathcal{C}\in\Cconf(G_k)$, $\mathcal{C}'\in\Cconf(G_{k'})$:
        $f(\mathcal{C},\mathcal{C}')\ge\delta_{\min}>0$.
  \item[(C3) No-beneficial-split:] Adding any extra group boundary (splitting a group
        into two) does not reduce the SSP-$f$ cost. Formally:
        for every feasible grouping $\calG$ with $K>\Kstar$ groups,
        $\OPT_f(\calG)\ge\OPT_f(\calG')$ for some JGP-optimal $\calG'$ with
        $\Kstar$ groups.
\end{description}
\end{theorem}
\end{thmbox}

\begin{proofbox}
\begin{proof}
\textbf{Necessity.}

$\lnot$(C1): If $f(\mathcal{C},\mathcal{C})>0$ for some $\mathcal{C}$, then
repeated uses of $\mathcal{C}$ incur positive cost, which is unphysical (no switch
occurs). More concretely: within a single group's processing, the magazine does not
change, so any well-defined switching cost must assign zero cost to
``no change''. If (C1) fails, $f$ models a cost even for non-events, and the
SSP-$f$ objective no longer corresponds to any meaningful switching count.
Any valid switching cost function must satisfy (C1).

$\lnot$(C2): Suppose $f(\mathcal{C},\mathcal{C}')=0$ for some
$\mathcal{C}\in\Cconf(G_k)$, $\mathcal{C}'\in\Cconf(G_{k'})$, $G_k\ne G_{k'}$.
Split $G_k$ into two sub-groups $G_k^{(1)},G_k^{(2)}$ with the same configuration
$\mathcal{C}$ for both, and insert $\mathcal{C}'$ between them.
The transition $\mathcal{C}\to\mathcal{C}'\to\mathcal{C}$ costs $0+0=0$.
This produces a $(\Kstar+1)$-group solution with the same SSP-$f$ cost as the
$\Kstar$-group solution, so the $(\Kstar+1)$-group solution is also optimal ---
collapse fails.

$\lnot$(C3): By definition, if some $K>\Kstar$ grouping achieves strictly lower
$f$-cost than all $\Kstar$-group solutions, the SSP-$f$ optimum uses $K>\Kstar$
groups, violating \cref{def:collapse}.

\textbf{Sufficiency.}
Under (C2), any solution with $K$ groups incurs at least $K-1$ inter-group transitions,
each costing $\ge\delta_{\min}>0$, giving total cost $\ge(K-1)\delta_{\min}$.
For $K>\Kstar$: $(K-1)\delta_{\min}>(\Kstar-1)\delta_{\min}$.
The $\Kstar$-group optimum $c^*_{\Kstar}$ satisfies $c^*_{\Kstar}\ge(\Kstar-1)\delta_{\min}$.
Under (C3), $c^*_{\Kstar}\le\OPT_f(K)$ for all $K>\Kstar$.
Hence every optimal solution uses exactly $\Kstar$ groups. \qed
\end{proof}
\end{proofbox}

\begin{remark}
Condition (C3) is the hardest to verify in general. The threshold and setup-matrix
variants below provide concrete sufficient conditions for (C3) in terms of instance
parameters.
\end{remark}

\subsection{Threshold SSP (New Variant)}
\label{sec:threshold}

\begin{definition}[Threshold SSP]
For $\theta\ge 0$:
$f_\theta(\mathcal{C},\mathcal{C}')=\max(0,|\mathcal{C}\setminus\mathcal{C}'|-\theta)$.
Switching $\le\theta$ tools is free; excess switches incur unit cost.
\end{definition}

\begin{thmbox}
\begin{theorem}[Threshold Collapse]
\label{thm:threshold}
For a JGP-optimal grouping $\calG^*$ with lookahead configurations, define
\[
  \theta^*(\calG^*)=\max_{k=1,\ldots,\Kstar-1}|T_{G_k}\setminus T_{G_{k+1}}|.
\]
For $\theta\ge\theta^*(\calG^*)$: SSP-$f_\theta$ achieves cost $0$ on $\calG^*$
and collapses to JGP. For $0\le\theta<\theta^*(\calG^*)$: SSP-$f_\theta$ has
strictly positive optimal cost and interpolates between SSP and TSIP.
\end{theorem}
\end{thmbox}

\begin{proof}
\textit{Collapse:}
Each lookahead transition costs $|T_{G_k}\setminus T_{G_{k+1}}|\le\theta^*\le\theta$,
so $f_\theta(\mathcal{C}_k^*,\mathcal{C}_{k+1}^*)=0$ for all $k$.
Total cost $=0=\OPT_{\mathrm{SSP}\text{-}f_\theta}$ (since $f_\theta\ge 0$).
Condition (C2) of \cref{thm:collapse} holds with $\delta_{\min}=1$ (any transition
between groups using different configurations costs at least $1-\theta\cdot 0=1$
after thresholding if $\theta=0$; more generally $\delta_{\min}=
\max(0,\min_{\mathrm{inter-group}}\sw{\cdot}{\cdot}-\theta)$ which is positive
when $\theta<\min_{\mathrm{inter-group}}\sw{\cdot}{\cdot}$, i.e.\ below the minimum
inter-group switching cost, which is $>0$ in any JGP-optimal grouping).
Condition (C3) holds because the $0$-cost $\Kstar$-group solution cannot be improved.

\textit{Non-collapse for $\theta<\theta^*$:}
There exists a consecutive pair $(G_{k^*},G_{k^*+1})$ with
$|T_{G_{k^*}}\setminus T_{G_{k^*+1}}|>\theta$.
By \cref{lem:transition}, the minimum transition cost for this pair over all valid
configurations is $w^*_{k^*}=|T_{G_{k^*}}\cup T_{G_{k^*+1}}|-C>0$.
After thresholding: $f_\theta\ge\max(0,w^*_{k^*}-\theta)>0$ for $\theta<w^*_{k^*}$.
Hence $\OPT_{\mathrm{SSP}\text{-}f_\theta}>0$. \qed
\end{proof}

\begin{remark}
The collapse threshold $\theta^*(\calG^*)$ depends on the grouping. The grouping with
highest inter-group tool overlap (selected by \cref{thm:overlap}) minimises $\theta^*$
and makes collapse easier to achieve in practice.
\end{remark}

\subsection{Setup-Matrix SSP (New Variant)}
\label{sec:setupmatrix}

\begin{definition}[Setup-Matrix SSP]
Let $W\in\mathbb{R}^{m\times m}_+$ with $W_{tt}=0$ and $W_{tt'}\ge 0$ for $t\ne t'$.
$f_W(\mathcal{C},\mathcal{C}')=\sum_{t\in\mathcal{C}\setminus\mathcal{C}'}
\min_{t'\in\mathcal{C}'\setminus\mathcal{C}}W_{tt'}$.
Here $W_{tt'}$ is the cost of removing $t$ when $t'$ is loaded in its place,
modelling sequence-dependent slot-reuse costs.
\end{definition}

\begin{thmbox}
\begin{theorem}[Setup-Matrix Collapse Condition]
\label{thm:setupmatrix}
If $W_{tt'}\ge\gamma>0$ for all $t\ne t'$, then for any consecutive group pair and
any valid configurations: $f_W(\mathcal{C}_k,\mathcal{C}_{k+1})\ge\gamma\cdot w^*_k$,
where $w^*_k=\max(0,|T_{G_k}\cup T_{G_{k+1}}|-C)$ from \cref{lem:transition}.
Condition (C2) holds with $\delta_{\min}=\gamma\cdot\min_k w^*_k>0$.
SSP-$f_W$ collapses to JGP when additionally (C3) holds:
$\gamma\cdot w^*_k$ exceeds the cost reduction from splitting any group.
\end{theorem}
\end{thmbox}

\begin{proof}
Each removed tool $t\in\mathcal{C}_k\setminus\mathcal{C}_{k+1}$ incurs cost
$\min_{t'\in\mathcal{C}_{k+1}\setminus\mathcal{C}_k}W_{tt'}\ge\gamma$
(the set $\mathcal{C}_{k+1}\setminus\mathcal{C}_k$ is non-empty and all $t'\ne t$).
Summing: $f_W\ge\gamma|\mathcal{C}_k\setminus\mathcal{C}_{k+1}|\ge\gamma\cdot w^*_k$.
For inter-group pairs, $w^*_k>0$ (adjacent groups cannot share a configuration
in a JGP-optimal grouping), so (C2) holds. \qed
\end{proof}

\subsection{Machine Learning for Collapse Variants}
\label{sec:ML-collapse}

\subsubsection{Learning the Collapse Threshold}

Train a regressor on instance features
$(n,m,C,\mathbb{E}[|T_j|],\mathrm{Var}[|T_j|],\max|T_j|,\rho)$
(where $\rho$ is tool co-occurrence density) to predict $\theta^*$.
At inference: predict $\hat\theta^*$; solve SSP-$f_{\hat\theta^*}$;
if accurate, the solution is JGP-optimal with zero threshold-cost.

\subsubsection{Meta-Learning a Collapse-Inducing Cost Function (New)}
\label{sec:metalearning}

Learn a parameterised $f_\phi=f_{W_\phi}$ (setup-matrix variant) satisfying
both collapse and fidelity to the original SSP:
\begin{equation}
  \label{eq:bilevel}
  \min_\phi\;\sum_I|\OPT_{\mathrm{SSP}\text{-}f_\phi}(I)-\OPT_{\mathrm{SSP}}(I)|
  +\lambda R(\phi),
\end{equation}
where
$R(\phi)=\sum_I\max(0,\OPT_{f_\phi}(I,K\le\Kstar)-\OPT_{f_\phi}(I,K=\Kstar))$
enforces collapse ($R(\phi)=0\iff$ SSP-$f_\phi$ collapses to JGP on training instances).

This is surrogate optimisation: SSP-$f_\phi$ (tractable via collapse) approximates
the original SSP (hard), with the two objectives in~\eqref{eq:bilevel} trading
off tractability against fidelity. Convergence guarantees and sample complexity
for this bilevel formulation are open.

\subsubsection{Lagrangian Parameter Prediction}

Rather than running the full subgradient algorithm of \cref{sec:lagrangian} on each
new instance, train a model $\hat\lambda^*(I)=f_\psi(\mathrm{features}(I))$ from
solved instances~\cite{Khalil2016}.
Predict $\hat\lambda^*$; solve $L(\hat\lambda^*)$ in one pass.
Analogous warm-starting approaches reduce solve times by 50--90\% in MIP settings.

% ─────────────────────────────────────────────────────────────────────────────
\section{Summary}
\label{sec:summary}
% ─────────────────────────────────────────────────────────────────────────────

\begin{table}[h!]
\centering
\caption{Key results: origin and status.}
\label{tab:summary}
\renewcommand{\arraystretch}{1.3}
\small
\begin{tabularx}{\textwidth}{@{}lXXl@{}}
\toprule
\textbf{Part} & \textbf{Result} & \textbf{Basis} & \textbf{Origin}\\
\midrule
I & Gap $H-\OPT=(K-1)s$; tight construction and full proof
  & Mandatory-exit LB (\cref{lem:LB}) + lookahead UB; both tight
  & \textbf{New}\\
I & Ratio $H/\OPT=(C-ov)/(r-ov)$; unbounded
  & Follows from construction
  & \textbf{New}\\
I & Gap $=0$ when $s=0$; SSP reduces to TSP on mandatory sets
  & Consistent with \cite{Ghiani2007}
  & Known\\
II & GTSP for fixed grouping
  & Specialisation of \cite{Ghiani2007} to fixed $\calG$
  & Framing new\\
II & Optimal ordering $=$ MWHP on overlap graph (\cref{thm:overlap})
  & New equivalence; tractable via Held--Karp
  & \textbf{New}\\
II & Lagrangian for JGP/SSP Pareto frontier
  & Standard Lagrangian duality \cite{Wolsey1998}; application is new
  & Application new\\
II & Column generation for SSP
  & \cite{daSilva2021}; reused here as grouping scorer
  & Reference\\
II & GNN grouping scorer; RL for joint grouping+sequencing
  & Architectures from \cite{Xu2019,Kool2019}; application new
  & \textbf{New}\\
II & Learning-augmented CG; contrastive ranking
  & Framework from \cite{Gasse2019}; application new
  & \textbf{New}\\
III & TSIP $\equiv$ JGP (baseline)
  & \cite{TangDenardo1988}
  & Known\\
III & General collapse characterisation (C1)+(C2)+(C3) (\cref{thm:collapse})
  & New necessary and sufficient conditions
  & \textbf{New}\\
III & Threshold SSP collapse (\cref{thm:threshold})
  & Corollary of \cref{thm:collapse}; new variant
  & \textbf{New}\\
III & Setup-Matrix SSP and $\gamma$-collapse condition (\cref{thm:setupmatrix})
  & New variant; new collapse proof
  & \textbf{New}\\
III & Meta-learning collapse-inducing $f_\phi$ (bilevel, \cref{sec:metalearning})
  & New open problem; novel direction
  & \textbf{New}\\
\bottomrule
\end{tabularx}
\end{table}

% ─────────────────────────────────────────────────────────────────────────────
\section{Open Research Questions}
\label{sec:open}
% ─────────────────────────────────────────────────────────────────────────────

\begin{enumerate}[label=\textbf{Q\arabic*.},leftmargin=*]
  \item \textbf{(I)} Is the MWHP score on $H_{\mathrm{ov}}$ a provably tight
        proxy for $\OPT(\calG)$ (not just under full lookahead)?
        Proving approximation ratios for groupings selected by this criterion is open.
  \item \textbf{(II)} What is the complexity of the CG pricing problem for SSP?
        Fixed-parameter tractability in $\Kstar$ would make branch-and-price practical.
  \item \textbf{(II)} Can GTSP cluster sizes be polynomially bounded for structured
        tool families (laminar, interval)? This would make \cref{thm:GTSP} directly tractable.
  \item \textbf{(III)} Does a function $f$ exist that (a) makes SSP-$f$ poly-time,
        (b) collapses to JGP, and (c) preserves $\OPT_{\mathrm{SSP}}$ within a constant?
        Likely no in general (implies P$=$NP), but open for parameterised versions.
  \item \textbf{(III/ML)} Convergence and sample complexity for the bilevel
        formulation~\eqref{eq:bilevel} are open; the landscape of $R(\phi)$ is
        non-convex and potentially discontinuous.
\end{enumerate}

% ─────────────────────────────────────────────────────────────────────────────
\begin{thebibliography}{99}

\bibitem{TangDenardo1988}
C.S.~Tang and E.V.~Denardo, ``Models arising from a flexible manufacturing machine,''
\textit{Operations Research}, vol.~36, no.~5, pp.~767--784, 1988.
[Original SSP; KTNS algorithm; TSIP $\equiv$ JGP.]

\bibitem{Crama1994}
Y.~Crama, O.E.~Flippo, J.~van de~Klundert, and F.C.R.~Spieksma,
``The tool switching problem revisited,''
\textit{European Journal of Operational Research}, vol.~74, pp.~331--342, 1994.
[NP-hardness; tool conflict graph.]

\bibitem{CramaKlundert1999}
Y.~Crama and J.~van de~Klundert,
``Worst-case performance of approximation algorithms for tool management problems,''
\textit{Naval Research Logistics}, vol.~46, pp.~445--462, 1999.

\bibitem{Laporte2004}
G.~Laporte, J.J.~Salazar-Gonz\'alez, and F.~Semet,
``Exact algorithms for the job sequencing and tool switching problem,''
\textit{IIE Transactions}, vol.~36, pp.~37--45, 2004.

\bibitem{Catanzaro2015}
D.~Catanzaro, L.~Gouveia, and M.~Labb\'e,
``Improved ILP formulations for the job sequencing and tool switching problem,''
\textit{EJOR}, vol.~244, pp.~766--777, 2015.

\bibitem{daSilva2021}
T.T.~da~Silva, A.A.~Chaves, L.A.N.~Lorena, and L.C.~Coelho,
``Multicommodity flow formulation and column generation for the SSP,'' 2021.
[LP lower bound $\OPT\ge m-C$; column generation.]

\bibitem{Ghiani2007}
G.~Ghiani, A.~Grieco, and R.~Musmanno,
``Reformulation and solution of the tool switching problem as a least-cost
Hamiltonian cycle problem,''
\textit{Journal of Intelligent Manufacturing}, vol.~18, pp.~523--531, 2007.
[SSP $=$ minimum-cost Hamiltonian cycle on configurations.]

\bibitem{NoonBean1993}
C.E.~Noon and J.C.~Bean,
``An efficient transformation of the generalized traveling salesman problem,''
\textit{INFORMS Journal on Computing}, vol.~5, pp.~571--575, 1993.

\bibitem{Wolsey1998}
L.A.~Wolsey, \textit{Integer Programming}. Wiley, 1998.
[Lagrangian relaxation and subgradient methods.]

\bibitem{Xu2019}
K.~Xu, W.~Hu, J.~Leskovec, and S.~Jegelka,
``How powerful are graph neural networks?'' \textit{ICLR 2019}.

\bibitem{Khalil2016}
E.~Khalil et al., ``Learning combinatorial optimization algorithms over graphs,''
\textit{NeurIPS 2016}.

\bibitem{Kool2019}
W.~Kool, H.~van Hoof, and M.~Welling,
``Attention, learn to solve routing problems!'' \textit{ICLR 2019}.

\bibitem{Gasse2019}
M.~Gasse et al., ``Exact combinatorial optimisation with GCNNs,'' \textit{NeurIPS 2019}.

\end{thebibliography}

\end{document}
